\section{Related Work}
\label{sec:related}

\paragraph{Syntax-Guided Synthesis.} 
The original formulation of inductive program synthesis is to search for a program in a hypothesis space (programming language) that is consistent with a specification (such as IO examples). However, this search is often intractable because of the large (potentially infinite) hypothesis space. One of the key ideas to make this search tractable is to provide the synthesizer a sketch of the desired program in addition to the examples, for example in \cite{Sketch:Main} and \cite{DBLP:conf/pldi/FeserCD15}. The program sketch in addition to providing structure to the search space also allows users to provide additional insights. This approach has been generalized in a framework called Syntax-Guided Synthesis (\textsc{Sygus})~\cite{sygus}. Our $\sygurl$ approach is inspired by \textsc{Sygus} in the sense that we also use a high-level grammar to constrain the shape of the possible learnt policies in a policy language grammar. However, unlike \textsc{Sygus} and previous sketch-based synthesis approaches that use logical constraints as specification, \sygurl searches for policies with quantitative objectives.
%\vspace*{-6pt}
\paragraph{Imitation Learning.} Imitation learning~\cite{schaal1999imitation} has been a successful paradigm for reducing the sample complexity of reinforcement learning algorithms by allowing the agent to leverage the additional supervision provided in terms of expert demonstrations for the desired behaviors. The \textsc{DAgger} (Dataset Aggregation) algorithm~\cite{dagger} is an iterative algorithm for imitation learning that learns stationary deterministic policies, where in each iteration $i$ it uses the current learnt policy $\pi_i$ to collect new trajectories and adds them to the dataset $D$ of all previously found trajectories. The policy for the next iteration $\pi_{i+1}$ is a policy that best mimics the expert policy $\pi^*$ on the whole dataset $D$. Our Neurally Directed Program Search (\algo) is inspired by the \textsc{DAgger} algorithm, where we use the trained DeepRL agent as the expert (oracle), and iteratively perform IO augmentation for unseen input states explored by our synthesized policy with the current best reward. However, one key difference is that \algo uses the expert trajectories to only guide the local program search in our policy language grammar to find a policy with highest rewards, unlike the imitation learning setting where the goal is to match the expert demonstrations perfectly.
%\vspace*{-6pt}
\paragraph{Neural Program Synthesis and Induction.} Many recent efforts use neural networks for learning programs. These efforts have two flavors. In \emph{neural program induction}, the goal is to learn a network that encodes the program semantics using internal weights. These architectures typically augment neural networks with differentiable computational substrates such as memory (Neural Turing Machines~\cite{ntm}), modules (Neural RAM~\cite{nram}) or data-structures such as stacks~\cite{stackrnn}, and formulate the program learning problem in an end-to-end differentiable manner. In \emph{neural program synthesis}, the architectures generate programs directly as outputs using multi-task transfer learning (e.g. \textsc{RobustFill}~\cite{RobustFill}, \textsc{DeepCoder}~\cite{DeepCoder}, \textsc{Bayou}~\cite{Bayou}), where the network weights are used to guide the program search in a DSL. There have also been some recent approaches to use \rl for learning to search programs in DSLs~\cite{rlkarel,rlbf}. Our approach falls in the category of program synthesis approaches where we synthesize policies in a policy language. However, we learn richer policy programs with continuous parameters using the \algo algorithm.
%\vspace*{-6pt}
\paragraph{Interpretable Machine Learning.} Many recent efforts in deep learning aim to make deep networks more interpretable \citep{NN:Understanding,NN:Mythos, SymbolicRL, SymbRL2, SymbRL3, SymbRL4}. There are three key approaches explored for interpreting DNNs: i) generate input prototypes in the input domain that are representatives of the learned concept in the abstract domain of the top-level of a DNN, ii) explaining DNN decisions by relevance propagation and computing corresponding representative concepts in the input domain, and iii) Using symbolic techniques to explain and interpret a DNN. 
%See for example, \cite{SymbolicRL}, \cite{SymbRL2}, \cite{SymbRL3}, \cite{SymbRL4}. 
Our work differs from these approaches in that we are replacing the DRL model with human readable source code, that is programmatically synthesized to mimic the policy found by the neural network. Working at this level of abstraction provides a method to apply existing synthesis techniques to the problem of making DRL models interpretable.
%\vspace*{-6pt}
\paragraph{Verification of Deep Neural Networks.} Reluplex~\cite{reluplex} is an SMT solver that supports linear real arithmetic with ReLU constraints, and has been used to verify several properties of DNN-based airborne collision avoidance systems, such as not producing erroneous alerts and uniformity of alert regions. Unlike Reluplex, our framework generates interpretable program source code as output, where we can use traditional symbolic program verification techniques~\cite{king1976symbolic} to prove program properties.
